\section{Definitions}\label{sec:definitions}

We begin with the finite-dimensional lattice model and the observables whose averages will converge.
We then specify the continuum correlations, their decomposition into phases and the relativistic fields reconstructed from one phase.
The last definitions give the two short-distance tests in the theorem.
Throughout, $\N=\{0,1,\ldots\}$; empty products equal $1$, and correlation functions of order zero are included.

\subsection{The lattice and its measure}\label{subsec:lattice-family}

\begin{definition}[Gauge group]\label{def:gauge}
A gauge group consists of an integer $m\ge1$, a closed subgroup $G$ of the group $U(m)$ of unitary $m\times m$ matrices, and a real linear subspace $\g$ of the complex matrix space $M_m(\C)$ such that
\begin{equation}\label{eq:gauge}
X\in\g\quad\Longleftrightarrow\quad \exp(sX)\in G\text{ for every }s\in\R\,.
\end{equation}
Its Lie algebra $\g$ is simple if $[X,Y]=XY-YX\ne0$ for some $X,Y\in\g$, and every real subspace $J\subset\g$ with $[\g,J]\subset J$ equals $0$ or $\g$.
\end{definition}

The matrices in $G$ will be the variables of the lattice model.
The specified embedding in $U(m)$ fixes the trace used in its energy.
We do not assume that $G$ is connected.

\begin{definition}[Periodic lattice]\label{def:lattice}
For $G$ as in Definition~\ref{def:gauge} and a positive integer $\ell$, let $\Lambda_\ell=(\Z/\ell\Z)^4$ and write $e_\mu$, $1\le\mu\le4$, for its coordinate unit vectors.
The sites are the points of $\Lambda_\ell$.
A positive link $(y,\mu)$ joins $y$ to $y+e_\mu$.
A link configuration is a map
\begin{equation}\label{eq:link-configuration}
U:\Lambda_\ell\times\{1,2,3,4\}\longrightarrow G\,.
\end{equation}
The value assigned to the reversed link is $U(y,\mu)^{-1}$.
\end{definition}

A link matrix describes parallel transport along one edge.
Multiplying four of these matrices around an elementary square compares transport along its two coordinate paths.
This product, the plaquette matrix, is the local quantity from which the Wilson energy is formed.

\begin{definition}[Plaquettes and energy]\label{def:plaquette}
For a configuration as in Definition~\ref{def:lattice}, set
\begin{equation}\label{eq:wilson-energy}
\begin{aligned}
U_{y;\mu\nu}&=U(y,\mu)U(y+e_\mu,\nu)U(y+e_\nu,\mu)^{-1}U(y,\nu)^{-1},\\
E_\ell(U)&=\sum_{y\in\Lambda_\ell}\sum_{\mu<\nu}
\left(1-\frac{1}{m}\operatorname{Re}\tr U_{y;\mu\nu}\right)\,.
\end{aligned}
\end{equation}
\end{definition}

The energy counts each unordered pair of coordinate directions once.
Each summand lies in $[0,2]$, so $0\le E_\ell\le12\ell^4$.
The probability measure weights configurations according to this energy.

\begin{definition}[Wilson measure]\label{def:wilson-law}
Let $G$ be as in Definition~\ref{def:gauge} and let $h$ be normalised Haar measure on $G$, the left-translation-invariant probability measure.
For $\beta\in\R$, the Wilson measure on the configurations of Definition~\ref{def:lattice} is
\begin{equation}\label{eq:wilson-law}
\begin{aligned}
Z_{\ell,\beta}&=\int e^{-\beta E_\ell(U)}\prod_{y,\mu}h(dU(y,\mu)),\\
\mu_{\ell,\beta}(dU)&=\frac{e^{-\beta E_\ell(U)}}{Z_{\ell,\beta}}\prod_{y,\mu}h(dU(y,\mu))\,.
\end{aligned}
\end{equation}
\end{definition}

The bound on $E_\ell$ makes $Z_{\ell,\beta}$ finite and strictly positive for every real $\beta$.
Increasing $\beta$ gives greater relative weight to configurations with smaller energy.
This parameter is called the bare inverse coupling: it is specified in the original lattice measure, before any variables are integrated out.
Changing coordinates in the internal matrix space independently at each site acts as follows.

\begin{definition}[Gauge transformations]\label{def:gauge-transform}
For configurations as in Definition~\ref{def:lattice} and $\gamma:\Lambda_\ell\to G$, define
\begin{equation}\label{eq:gauge-transform}
U^\gamma(y,\mu)=\gamma(y)U(y,\mu)\gamma(y+e_\mu)^{-1}\,.
\end{equation}
An observable $F$ on the configurations of Definition~\ref{def:lattice} is gauge invariant if $F(U^\gamma)=F(U)$ for every $U$ and $\gamma$.
\end{definition}

\subsection{Local observables and the continuum limit}\label{subsec:observables}

An observable is a function of the link matrices; its expectation is a measurement in the finite lattice model.
To pass to a field at a point, we use observables that depend on links near one site.
Their coefficients may change as the lattice is refined.
To follow one observable through these changes, we give it a fixed index, separate from its formula on a particular lattice.

\begin{definition}[Observable indices and insertions]\label{def:observables}
Let $\mathcal T$ be a set of observable indices.
An insertion with index $v\in\mathcal T$ at a site $y\in\Lambda_\ell$ is a function on link configurations (Definition~\ref{def:lattice}):
\begin{equation}\label{eq:observable-insertion}
I(v,y;\cdot):G^{\Lambda_\ell\times\{1,2,3,4\}}\longrightarrow\C,
\qquad U\longmapsto I(v,y;U)\,.
\end{equation}
For a finite set $\Sigma\subset\Z^4$, write $\Sigma\bmod\ell$ for its image in $\Lambda_\ell$.
The insertion is supported on the links inside $y+(\Sigma\bmod\ell)$ if it takes the same value on any two configurations that agree on all links whose two endpoints belong to this set.
\end{definition}

The index set is arbitrary at this stage; no linear structure or degree is required.
Our construction uses pairs consisting of a degree and a gauge-invariant polynomial in the Taylor coefficients of a classical connection, a smooth $\g$-valued one-form describing parallel transport (Definition~\ref{def:uv-labels}).
That polynomial specifies the classical observable to be represented, while the insertion is its actual function of link matrices.
The continuum field obtained from an observable index will be part of the hierarchy below.
We next choose the physical sizes of the lattices.
The index $N$ refines the spacing and $p$ increases the volume.

\begin{definition}[Wilson family]\label{def:family}
For a gauge group (Definition~\ref{def:gauge}), a Wilson family $\W$ consists of strictly increasing positive integer sequences $(D_N)_{N\in\N}$ and $(L_p)_{p\in\N}$, real numbers $\beta_{N,p}$, normalised Haar measure $h$ on $G$, a set $\mathcal T$ of observable indices, and insertions $I_{N,p}(v,y;U)$ as in Definition~\ref{def:observables} on $\Lambda_{L_pD_N}$.
We write
\begin{equation}\label{eq:physical-scales}
a_N=\frac{1}{D_N},\qquad \ell_{N,p}=L_pD_N,\qquad
\mu_{N,p}=\mu_{\ell_{N,p},\beta_{N,p}}\,.
\end{equation}
\end{definition}

The lattice has $L_pD_N$ sites per side, separated by physical distance $a_N$; thus its physical side is $L_p$.
Taking $N\to\infty$ at fixed $p$ refines the lattice without enlarging the box.
Taking $p\to\infty$ enlarges the box.
The spacing $a_N$ is a short-distance cutoff; its limit to zero is called the ultraviolet limit.
The family also specifies the observables whose averages are to have limits.

\begin{definition}[Regularity]\label{def:regular}
A Wilson family (Definition~\ref{def:family}) is regular if:
\begin{enumerate}[label=\textup{Reg\arabic*.},ref=Reg\arabic*]
\item\label{reg:1} There exists $b_*>0$ such that, for every $\varepsilon>0$, there exists $N_0$ with
\begin{equation}\label{eq:running}
\left|\frac{\beta_{N,p}}{\log\left(\frac{1}{a_N}\right)}-b_*\right|<\varepsilon
\qquad(N\ge N_0,\ p\in\N)\,.
\end{equation}
\item\label{reg:2} Each $U\mapsto I_{N,p}(v,y;U)$ is continuous.
\item\label{reg:3} Each insertion is gauge invariant in Definition~\ref{def:gauge-transform}.
\item\label{reg:4} For every $v\in\mathcal T$, there exists a finite set $\Sigma_v\subset\Z^4$ such that every $I_{N,p}(v,y;\cdot)$ is supported on the links inside $y+(\Sigma_v\bmod\ell_{N,p})$ in Definition~\ref{def:observables}.
\end{enumerate}
\end{definition}

The same finite set in~\ref{reg:4} works at every spacing, volume and site.
Its diameter in physical units tends to zero with $a_N$.
Condition~\ref{reg:1} is a logarithmic growth law for the bare inverse coupling, uniform in the volume; its positive slope is left unspecified here.

\begin{definition}[Physical positions and moments]\label{def:moments}
Let $r_\ell(j)$ be the representative of $j\in\Z/\ell\Z$ in
$\left(-\frac{\ell}{2},\frac{\ell}{2}\right]$.
The physical position of $y\in\Lambda_{\ell_{N,p}}$ is
\begin{equation}\label{eq:physical-position}
\hat y=a_N\bigl(r_{\ell_{N,p}}(y^\mu)\bigr)_{\mu=1}^4\,.
\end{equation}
For a Wilson family (Definition~\ref{def:family}) and $v=(v_1,\ldots,v_n)\in\mathcal T^n$, define the moment and its smearing by a function $f$ by
\begin{equation}\label{eq:smear}
\begin{aligned}
M_{N,p}(v;y)&=\int\prod_{i=1}^n I_{N,p}(v_i,y_i;U)\,\mu_{N,p}(dU),\\
C_{N,p}(v;f)&=a_N^{4n}\sum_{y\in\Lambda_{\ell_{N,p}}^n}
 f(\hat y_1,\ldots,\hat y_n)M_{N,p}(v;y)\,.
\end{aligned}
\end{equation}
\end{definition}

A correlation is an average of a product of observables.
Smearing averages its spatial arguments against a test function, so the lattice sum can be compared with a continuum integral.
We test only configurations of distinct points.

\begin{definition}[Separated tests]\label{def:tests}\label{def:separated-test}
For $n\in\N$, put
\begin{equation}\label{eq:conf}
\Conf n=\{x\in(\R^4)^n:x_i\ne x_j\text{ whenever }i\ne j\}\,.
\end{equation}
A separated test is a smooth function $f:(\R^4)^n\to\C$ whose closed support is compact and contained in $\Conf n$.
We write $\supp f$ for the closed support, the closure of the set on which $f$ is nonzero.
\end{definition}

Such a test vanishes near every coincidence of two points.
It can depend jointly on all $n$ points and need not be a product of one-point tests.
To compare a torus with Euclidean space, its support must also stay inside the physical box.

\begin{definition}[Tests supported inside a volume]\label{def:admissible}\label{def:admissible-test}
A separated test $f$ (Definition~\ref{def:tests}) is supported inside volume $p$ if
\begin{equation}\label{eq:admissible}
\supp f\subset
\left(\left(-\frac{L_p}{2},\frac{L_p}{2}\right)^4\right)^n\,.
\end{equation}
\end{definition}

\begin{definition}[Ordered limit]\label{def:ordered}
For a Wilson family (Definition~\ref{def:family}), $P:\N\to\N$ and kernels $K_n(v;x)$, convergence along $P$ means the following.
For every $n\in\N$, $v\in\mathcal T^n$ and separated test $f$ (Definition~\ref{def:tests}), the test is supported inside volume $P(k)$ eventually in $k$, and there exist $R_p\in\C$ such that
\begin{equation}\label{eq:ordered}
\begin{aligned}
\lim_{N\to\infty}C_{N,p}(v;f)&=R_p
&&\text{($f$ supported in volume $p$)},\\
\lim_{k\to\infty}R_{P(k)}&=\int_{(\R^4)^n}f(x)K_n(v;x)\,dx\,.
\end{aligned}
\end{equation}
\end{definition}

Here the moments are those of Definition~\ref{def:moments} and the support condition is Definition~\ref{def:admissible}.
The first limit uses the entire spacing sequence at every volume containing the test.
The second limit enlarges the volume along $P$.
At volumes not containing the support, $R_p$ is unconstrained.
A strictly increasing $P$ contains each fixed compact test support eventually, since $L_{P(k)}\to\infty$.

\begin{definition}[Periodic separation]\label{def:periodic}
For a Wilson family (Definition~\ref{def:family}) and volume $p$, the periodically separated configurations are
\begin{equation}\label{eq:periodic-separated}
\{x\in(\R^4)^n:x_i-x_j\notin L_p\Z^4\text{ whenever }i\ne j\}\,.
\end{equation}
\end{definition}

Distinct points can represent the same point of a torus if their difference is a period.
This definition excludes those coincidences as well, and is used for fixed-volume correlation kernels.

\subsection{The continuum objects}\label{subsec:continuum-objects}

\begin{definition}[Spacetime and Schwartz tests]\label{def:spacetime}
Write $x=(\vec x,t(x))\in\R^3\times\R$ for spacetime, with $t(x)=x^4$, Euclidean norm $|x|$ and
\begin{equation}\label{eq:spacetime}
\mathsf m(x)=t(x)^2-|\vec x|^2,
\qquad \overline V_+=\{x:|\vec x|\le t(x)\}\,.
\end{equation}
A vector is spacelike if $\mathsf m(x)<0$.
Let $\mathcal S=\Sch$, the smooth complex functions whose derivatives decay faster than every inverse power of $|x|$, with the corresponding Schwartz topology.
\end{definition}

Hilbert inner products are antilinear in the first variable.
The next definition lists the Hilbert space and fields before imposing their physical properties.
A field is smeared against a Schwartz function because its value at a single point need not be an operator.

\begin{definition}[Relativistic data]\label{def:relativistic}
A relativistic datum consists of a complex Hilbert space $\mathcal H$, a complex vector space $\mathcal L$ of field labels, a linear subspace $\mathcal D\subset\mathcal H$, a vector $\Omega\in\mathcal D$ called the vacuum, and fields $\Phi(\ell,f):\mathcal D\to\mathcal D$ linear over $\C$ in $\ell\in\mathcal L$ and $f\in\mathcal S$ (Definition~\ref{def:spacetime}).
The field label $\ell$ specifies which quantum field is smeared; the test function $f$ specifies how it is tested in space and time.
It includes an antilinear map $\ell\mapsto\ell^*$, a strongly continuous unitary translation representation $T(b)$ and a projection-valued Borel measure $E$ on $\R^4$.
Its diagonal finite measures satisfy $\mu_\psi(A)=\langle\psi,E(A)\psi\rangle$ on Borel sets.

There is also a strongly continuous unitary representation $U$ of a topological group $\mathcal P$, identified by a homeomorphism $\kappa$ with the proper orthochronous Poincar\'e group.
Writing $\kappa(g)=(\Lambda_g,b_g)$, this identification respects identity and the multiplication
\begin{equation}\label{eq:poincare}
(\Lambda,b)(\Lambda',b')=(\Lambda\Lambda',\Lambda b'+b)\,.
\end{equation}
Here $\Lambda$ is invertible, has determinant $1$, preserves $\mathsf m$ and maps $\overline V_+$ onto itself.
Pure translations $\tau(b)=\kappa^{-1}(1,b)$ satisfy $U(\tau(b))=T(b)$.
For each $g\in\mathcal P$ there is a complex linear bijection $\rho(g)$ of $\mathcal L$.
The action on tests is $(g\cdot f)(x)=f(\Lambda_g^{-1}(x-b_g))$.
\end{definition}

A projection-valued Borel measure here is normalised by $E(\R^4)=1$, has $E(\varnothing)=0$, satisfies $E(A)E(B)=E(A\cap B)$ and $E(A)^*=E(A)$ on Borel sets, and is countably additive in the strong operator topology.
One can assign arbitrary operators to non-Borel sets, which none of the assertions uses.
The map ${}^*$ on field labels is not required to be an involution.
No continuity or finite-dimensionality condition on $\rho$ is included in the data.

\begin{definition}[Wightman conditions]\label{def:wightman}
The relativistic datum of Definition~\ref{def:relativistic} satisfies the following conditions, with spacetime and test functions as in Definition~\ref{def:spacetime} and $p\cdot b$ the Euclidean pairing.
\begin{enumerate}[label=\textup{W\arabic*.},ref=W\arabic*,start=0]
\item\label{wightman:0} For every $\psi\in\mathcal H$ and $b\in\R^4$,
\begin{equation}\label{eq:spectrum}
\langle\psi,T(b)\psi\rangle=\int e^{ip\cdot b}\,\mu_\psi(dp),
\qquad E(\R^4\setminus\overline V_+)=0\,.
\end{equation}
Every vector fixed by all $T(b)$ belongs to $\C\Omega$.
\item\label{wightman:1} The vacuum has norm one, $\mathcal D$ is dense, and
\begin{equation}\label{eq:domain}
\mathcal D=\Span_{\C}\{\Phi(\ell_1,f_1)\cdots\Phi(\ell_j,f_j)\Omega:j\in\N,\ \ell_i\in\mathcal L,\ f_i\in\mathcal S\}\,.
\end{equation}
For every $\ell\in\mathcal L$ and $\psi\in\mathcal D$, the map $f\mapsto\Phi(\ell,f)\psi$ is continuous from $\mathcal S$ into $\mathcal H$ in norm.
For every $\ell,f$ and $\psi,\chi\in\mathcal D$,
\begin{equation}\label{eq:adjoint}
\langle\Phi(\ell,f)\psi,\chi\rangle=\langle\psi,\Phi(\ell^*,\bar f)\chi\rangle\,.
\end{equation}
\item\label{wightman:2} The maps $\rho$ form a representation: $\rho(1)=1$ and $\rho(gg')=\rho(g)\rho(g')$.
For every $g\in\mathcal P$, $U(g)\mathcal D\subset\mathcal D$ and $U(g)\Omega=\Omega$.
For every $g,\ell,f$ and $\psi\in\mathcal D$,
\begin{equation}\label{eq:covariance}
U(g)\Phi(\ell,f)\psi=\Phi(\rho(g)\ell,g\cdot f)U(g)\psi\,.
\end{equation}
\item\label{wightman:3} If $\mathsf m(x-y)<0$ for every $x\in\supp f$ and $y\in\supp f'$, then, for every $\ell,\ell'\in\mathcal L$ and $\psi\in\mathcal D$,
\begin{equation}\label{eq:locality}
\Phi(\ell,f)\Phi(\ell',f')\psi=\Phi(\ell',f')\Phi(\ell,f)\psi\,.
\end{equation}
\end{enumerate}
\end{definition}

The finite products of smeared fields in~\eqref{eq:domain} are called field words.
Their values on the vacuum span the domain itself; the empty word gives the vacuum.
The symmetry group is the Poincar\'e group itself rather than its double cover.
The sign convention in~\eqref{eq:spectrum} is converted to the usual Minkowski pairing by reversing spatial momentum; the forward cone is unchanged.

We collect the continuum objects in a hierarchy: correlations, data for their decomposition into phases, relativistic fields, and coefficients for local product expansions.

\begin{definition}[Hierarchy]\label{def:hierarchy}
A hierarchy over the Wilson family of Definition~\ref{def:family} consists of the following data.
\begin{enumerate}[label=\textup{H\arabic*.},ref=H\arabic*]
\item\label{hierarchy:1} A relativistic datum of Definition~\ref{def:relativistic}.
\item\label{hierarchy:2} For every $d\in\N$, a finite-dimensional complex normed space $F_d$ and a complex linear map $\iota_d:F_d\to\mathcal L$.
Set $\mathbb F=\bigsqcup_{d\in\N}F_d$, $\deg(d,A)=d$ and $\iota(d,A)=\iota_d A$.
\item\label{hierarchy:3} A map $\pi:\mathcal T\to\mathbb F$, assigning a continuum field to each observable index (Definition~\ref{def:observables}).
\item\label{hierarchy:4} A compact metric space $X$, called the phase space, with a Borel measure $\nu$ and a selected point $\xi_*\in X$.
\item\label{hierarchy:5} For every $n\in\N$ and $s\in X\sqcup\{\mix\}$, a correlation function $S^s(a;x)\in\C$ on $\mathbb F^n\times(\R^4)^n$.
The extra symbol $\mix$ distinguishes the average state supplied by the lattice limit from the phases $\xi\in X$.
\item\label{hierarchy:6} For every $n\in\N$, a continuation $\mathsf F_n(a;z,\vec x)\in\C$ on $\mathbb F^n\times\C^n\times(\R^3)^n$.
\item\label{hierarchy:7} For every $j\in\N$, a finite index set $\Ch j$ and fields $O_{j,\kappa}\in F_j$.
For every finite index set $I$, $a\in\mathbb F^I$ and $\kappa\in\Ch j$, there is a coefficient function $c_{j,\kappa}(a;x,o)$ on $(\R^4)^I\times\R^4$.
\item\label{hierarchy:8} An element $\Theta\in F_4$.
\end{enumerate}
\end{definition}

The maps in~\ref{hierarchy:2}--\ref{hierarchy:3} relate the observable indices, continuum fields and quantum field labels:
\begin{equation}\label{eq:observable-field-maps}
v\in\mathcal T\ \xmapsto{\pi}\ (d,A)\in\mathbb F
\ \xmapsto{\iota}\ \ell=\iota_d A\in\mathcal L\,.
\end{equation}
The field label specifies the operator family $f\mapsto\Phi(\ell,f)$ (Definition~\ref{def:relativistic}).
In our construction, the polynomial describes the classical observable and its index follows the lattice representative.

The degree is an index used in the product estimates below.
The definition imposes no inclusions between different $F_d$, no prescribed curvature interpretation and no compatibility of the $\iota_d$ across degrees.
The functions $S^s$ are defined even where points coincide, but the conclusions constrain them only at separated points.

\subsection{Reconstruction and phases}\label{subsec:reconstruction}

The mixture will be the limit of lattice moments.
The point $\xi_*$ instead selects the Euclidean correlations that continue to the relativistic fields.
These two roles are related by a phase average, not by identifying the two states.

\begin{definition}[Reconstruction properties]\label{def:reconstructed}
A hierarchy of Definition~\ref{def:hierarchy} has the reconstruction properties if the following hold.
Here $s\in X\sqcup\{\mix\}$, $n\in\N$ and $a\in\mathbb F^n$.
\begin{enumerate}[label=\textup{R\arabic*.},ref=R\arabic*]
\item\label{recon:1} Its relativistic datum satisfies \ref{wightman:0}--\ref{wightman:3}.
\item\label{recon:2} For every nonzero $\ell\in\mathcal L$, there exists $f\in\mathcal S$ with $\Phi(\ell,f)\ne0$.
This is faithfulness: no nonzero field label represents the zero field.
\item\label{recon:3} Each $\iota_d$ is injective.
\item\label{recon:4} Some $a\in\mathbb F$ and $f\in\mathcal S$ satisfy $\Phi(\iota(a),f)\ne0$.
\item\label{recon:5} The map $a\mapsto\iota(a)$ from $\mathbb F$ onto $\mathcal L$ is surjective.
\item\label{recon:6} The map $\pi:\mathcal T\to\mathbb F$ is surjective.
\item\label{recon:7} The measure $\nu$ has total mass one.
\item\label{recon:8} Every nonempty open subset of $X$ has positive $\nu$-measure; this is full support.
\item\label{recon:9} For every state $s$, order $n$, list $a\in\mathbb F^n$, $x\in\Conf n$ and permutation $\sigma$,
\begin{equation}\label{eq:permutation}
S^s((a_{\sigma(i)})_i;(x_{\sigma(i)})_i)=S^s(a;x)\,.
\end{equation}
\item\label{recon:10} For every $s,n,a$, the function $S^s(a;\cdot)$ is smooth on $\Conf n$.
\item\label{recon:11} For every $n,a$ and $x\in\Conf n$, the function $\xi\mapsto S^\xi(a;x)$ is continuous.
\item\label{recon:12} For every $n,a$ and $x\in\Conf n$,
\begin{equation}\label{eq:mixture}
\int_X S^\xi(a;x)\,\nu(d\xi)=S^\mix(a;x)\,.
\end{equation}
\item\label{recon:13} For every $n,a$ and fixed real spatial configuration $\vec x$, the function $z\mapsto\mathsf F_n(a;z,\vec x)$ is holomorphic on
\begin{equation}\label{eq:tube}
\mathcal Z_n=\{z\in\C^n:\operatorname{Im}z_1<\cdots<\operatorname{Im}z_n\}\,.
\end{equation}
\item\label{recon:14} For every $n,a$ and real configuration $x$ with $t(x_1)<\cdots<t(x_n)$,
\begin{equation}\label{eq:euclidean-root}
\mathsf F_n(a;(it(x_i))_i,(\vec x_i)_i)=S^{\xi_*}(a;x)\,.
\end{equation}
\item\label{recon:15} For every $n,a$, real heights $h_1<\cdots<h_n$ and tests $f_1,\ldots,f_n\in\mathcal S$,
\begin{multline}\label{eq:boundary}
\lim_{r\downarrow0}\int_{(\R^4)^n}\mathsf F_n(a;(t(x_i)+irh_i)_i,(\vec x_i)_i)\prod_{i=1}^n f_i(x_i)\,dx\\
=\langle\Omega,\Phi(\iota(a_1),f_1)\cdots\Phi(\iota(a_n),f_n)\Omega\rangle\,.
\end{multline}
\end{enumerate}
\end{definition}

The scalar on the right of~\eqref{eq:boundary} is the vacuum expectation of the indicated product of smeared fields.
In~\eqref{eq:boundary} we retain the formal source's convention that the integral of a nonintegrable function is zero.
No separate positive-height integrability assertion is part of \ref{recon:15}.
The integrals in the lattice limit and~\eqref{eq:mixture} are ordinary absolutely convergent integrals, by compact support and smoothness in the first case and by phase continuity and compactness in the second.

\subsection{Products and the trace test}\label{subsec:products}

The product expansion replaces a cluster by finitely many single fields at its centre.
Fields whose positions remain away from the cluster are called spectators.
One choice of coefficients must work with every list of such fields and in every state.
The threshold on the cluster radius, however, may depend on the fixed spectator fields.

\begin{definition}[Cluster geometry]\label{def:cluster}
For a hierarchy (Definition~\ref{def:hierarchy}) and finite index sets $I,J$, write $S^s(a,x;b,y)$ for the correlation on the concatenated lists, enumerated in one fixed order on $I\sqcup J$.
For $r>0$, $M\ge0$ and $\eta>0$, call $(x,o)$ an $(r,M,\eta)$-cluster if
\begin{equation}\label{eq:cluster}
|x_i-o|\le Mr\quad(i\in I),
\qquad |x_i-x_{i'}|\ge\eta r\quad(i\ne i')\,.
\end{equation}
For $c>0$, call $y$ $c$-clear of $o$ if
\begin{equation}\label{eq:clear}
|o-y_j|\ge c\quad(j\in J),
\qquad |y_j-y_{j'}|\ge c\quad(j\ne j')\,.
\end{equation}
Set $w(a)=\sum_{i\in I}\deg a_i$.
\end{definition}

\begin{definition}[Operator product expansion]\label{def:ope}
For a hierarchy (Definition~\ref{def:hierarchy}), the OPE property consists of the following estimates on clusters and distant points as in Definition~\ref{def:cluster}.
\begin{enumerate}[label=\textup{O\arabic*.},ref=O\arabic*]
\item\label{ope:1} For every finite $I$, $a\in\mathbb F^I$, $M\ge0$, $\eta,\varepsilon>0$, $j\in\N$ and $\kappa\in\Ch j$, there exists $\delta>0$ such that, for every $0<r<\delta$ and every $(r,M,\eta)$-cluster,
\begin{equation}\label{eq:ope-coefficient}
|c_{j,\kappa}(a;x,o)|\le r^{j-w(a)-\varepsilon}\,.
\end{equation}
\item\label{ope:2} For every finite $I,J$, $a\in\mathbb F^I$, $b\in\mathbb F^J$, $M\ge0$, $\eta,c,\varepsilon>0$ and $D\in\N$, there exists $\delta>0$ such that, for every $0<r<\delta$, every $(r,M,\eta)$-cluster, every $c$-clear spectator configuration and every $s\in X\sqcup\{\mix\}$,
\begin{equation}\label{eq:ope-remainder}
\left|S^s(a,x;b,y)-\sum_{j=0}^D\sum_{\kappa\in\Ch j}
 c_{j,\kappa}(a;x,o)S^s\bigl((j,O_{j,\kappa}),o;b,y\bigr)\right|\le r^{D+1-w(a)-\varepsilon}\,.
\end{equation}
\end{enumerate}
\end{definition}

The coefficient functions are those already chosen in \ref{hierarchy:7}.
They do not depend on spectators, state or truncation degree.
Both estimates have constant $1$, with the dependence absorbed into $\delta$.
For small enough $r$, $Mr<c$ ensures that all points in the tested correlations are separated.
These are finite-truncation estimates; they assert neither an infinite convergent expansion nor uniqueness of its coefficients.

\begin{definition}[The trace two-point function]\label{def:trace-pair}
For the field $\Theta\in F_4$ of Definition~\ref{def:hierarchy} and $r\in\R$, set
\begin{equation}\label{eq:trace-pair}
G_\Theta(r)=S^{\xi_*}\Bigl((4,\Theta),(4,\Theta);-\frac{r}{2}e_4,\frac{r}{2}e_4\Bigr)\,.
\end{equation}
\end{definition}

Free-field two-point functions are built from the energy $E_{\mathsf m}(p)$ of a particle with mass $\mathsf m$ and spatial momentum $p$.
The following class allows finitely many masses and derivatives, and finite sums of products of the resulting kernels.
We use these kernels only at positive separations $r>0$.
It specifies exactly the comparisons excluded by the theorem.

\begin{definition}[Finite Wick profiles]\label{def:wick}
A mode consists of $\mathsf m\ge0$, integers $j,q\in\N$ and a continuous real function $h$ on $\R^3$ satisfying $h(\lambda p)=\lambda^jh(p)$ for every $\lambda>0$.
Set
\begin{equation}\label{eq:mode}
E_{\mathsf m}(p)=\sqrt{\mathsf m^2+|p|^2},
\qquad K_{\mathsf m,q,h}(r)=\int_{\R^3}\frac{E_{\mathsf m}(p)^qh(p)}{2E_{\mathsf m}(p)}e^{-rE_{\mathsf m}(p)}\,dp\,.
\end{equation}
The integrand can be assigned value zero at a zero denominator, a set of measure zero.
A finite Wick profile is
\begin{equation}\label{eq:wick}
W(r)=\sum_{\alpha=1}^A c_\alpha\prod_{\beta=1}^{n_\alpha}
 K_{\mathsf m_{\alpha\beta},q_{\alpha\beta},h_{\alpha\beta}}\bigl(r\bigr),
\qquad c_\alpha\in\C,\quad A,n_\alpha\in\N\,.
\end{equation}
Empty products are allowed.
\end{definition}

\begin{definition}[Trace test]\label{def:trace}
The hierarchy of Definition~\ref{def:hierarchy} passes the trace test if the following hold, with $G_\Theta$ as in Definition~\ref{def:trace-pair} and Wick profiles as in Definition~\ref{def:wick}.
\begin{enumerate}[label=\textup{T\arabic*.},ref=T\arabic*]
\item\label{trace:1} For every $f\in\mathcal S$, $\langle\Omega,\Phi(\iota_4\Theta,f)\Omega\rangle=0$.
\item\label{trace:2} The number $A_G=\frac{3\dim_\R\g}{\pi^4}$ is positive.
\item\label{trace:3} The following limit holds in $\C$ over all positive radii:
\begin{equation}\label{eq:trace-law}
\lim_{r\downarrow0}r^8\left(\log\frac{1}{r}\right)^2G_\Theta(r)=\frac{3\dim_\R\g}{\pi^4}\,.
\end{equation}
\item\label{trace:4} For every finite Wick profile $W$ and every $\delta>0$, there exists $r\in(0,\delta)$ with $G_\Theta(r)\ne W(r)$.
\end{enumerate}
\end{definition}

The public trace test concerns the chosen degree-four field and the time axis.
Its definition alone supplies no stress tensor.
In the construction, a stress tensor is a tensor of fields whose integral identities express translation and rotation symmetry; its Euclidean trace is the sum of its diagonal components.
The construction identifies this field with a centred stress trace, as used in the proof of Proposition~\ref{prop:trace} below.

